\documentclass[10pt,a4paper]{article}
\usepackage[T1]{fontenc}
\usepackage[utf8]{inputenc}
\usepackage[czech]{babel}
\usepackage{lmodern}
\usepackage{microtype}
\usepackage{geometry}
\geometry{left=14mm,right=14mm,top=18mm,bottom=18mm}
\usepackage[table]{xcolor}
\usepackage{array,longtable,tabularx,booktabs}
\usepackage{fancyhdr,lastpage}
\usepackage{ragged2e}
\usepackage{hyperref}
\hypersetup{hidelinks,pdftitle={Rozpis výuky matematiky – 1. ročník SPŠ, 136 hodin podle tematických plánů a skutečné praxe}}
\setlength{\parindent}{0pt}
\setlength{\parskip}{5pt}
\setlength{\tabcolsep}{4pt}
\renewcommand{\arraystretch}{1.08}
\setlength{\headheight}{18pt}
\pagestyle{fancy}\fancyhf{}
\fancyhead[L]{\footnotesize Matematika -- rozpis výuky}
\fancyhead[R]{\footnotesize 1. ročník SPŠ}
\fancyfoot[L]{\footnotesize Rozpis hodin podle tematického plánu}
\fancyfoot[R]{\footnotesize \thepage/\pageref{LastPage}}
\renewcommand{\headrulewidth}{0.35pt}
\newcommand{\testr}[1]{\textbf{#1}}
\begin{document}
\begin{center}
{\LARGE\bfseries Rozpis výuky matematiky}\\[3mm]
{\large 1. ročník střední průmyslové školy}\\[3mm]
{\small Přehled plánované výuky podle tematických plánů SP\v{S} a reálných zápisů výuky}
\end{center}
\vspace{3mm}
\begin{tabularx}{\textwidth}{@{}Xrrrr@{}}\toprule
\textbf{Ročník} & \textbf{Hodin} & \textbf{Krátké testy} & \textbf{Souhrnné} & \textbf{Testů celkem}\\\midrule
1. ročník & 136 & 8 & 4 & 12\\\midrule
\textbf{Celkem} & \textbf{136} & \textbf{8} & \textbf{4} & \textbf{12}\\\bottomrule
\end{tabularx}
\vspace{4mm}
\section*{Charakter dokumentu}
\begin{itemize}\setlength{\itemsep}{2pt}
\item Dokument obsahuje rozpis výuky matematiky pro 1. ročník SP\v{S} v celkové dotaci \textbf{136 hodin}.
\item Rozpis je členěn po měsících a čtvrtletích. Každý očíslovaný řádek představuje \textbf{jednu plánovanou vyučovací hodinu}.
\item Zkratka \textbf{TP} odkazuje na odpovídající tematický okruh tematického plánu. U testů je označení uvedeno přímo v náplni hodiny.
\item V každém čtvrtletí jsou zařazeny \textbf{tři testy}: dva průběžné testy (zpravidla 15--20 minut) a jedna souhrnná písemná práce.
\end{itemize}

\section*{Z čeho rozpis vychází}
\begin{itemize}\setlength{\itemsep}{2pt}
\item Hodinová dotace vychází z tematického plánu SP\v{S} pro 1. ročník: \textbf{136 hodin}.
\item Obsah a pořadí témat zohledňují reálné zápisy výuky v dostupných třídách, zejména v 1.A (2024/25), 2.A (2025/26), 3.A (2025/26) a 3.D (2025/26).
\item Historické zápisy slouží jako podklad pro reálné pořadí témat a odhad časové náročnosti; nejde o doslovný přepis třídních knih.
\end{itemize}

\section*{Podkladové dokumenty}
\begin{itemize}\setlength{\itemsep}{2pt}
\item tematické plány SP\v{S},
\item zápisy výuky.
\end{itemize}

\section*{Odlišnosti oproti tematickým plánům}
\begin{itemize}\setlength{\itemsep}{2pt}
\item \textbf{1. ročník:} výrazněji je rozpracován blok výrokové logiky a množin; pořadí témat odpovídá reálné výuce a zachovává již vytvořené první přípravy.
\end{itemize}
\clearpage

\section*{1. ročník – logika, množiny, algebra a funkce}
\textit{Dotace 136 hodin. Čtvrtletní celky a tři testy v každém čtvrtletí jsou součástí hodinové dotace.}
\vspace{2mm}
\footnotesize
\begin{longtable}{@{}>{\RaggedLeft\arraybackslash}p{9mm}>{\RaggedRight\arraybackslash}p{15mm}>{\RaggedRight\arraybackslash}p{150mm}@{}}
\toprule \textbf{Hod.} & \textbf{TP} & \textbf{Náplň vyučovací hodiny} \\\midrule\endfirsthead
\toprule \textbf{Hod.} & \textbf{TP} & \textbf{Náplň vyučovací hodiny} \\\midrule\endhead
\bottomrule\endfoot
\addlinespace[1mm]\multicolumn{3}{@{}l}{\footnotesize\textbf{1. čtvrtletí} -- 2 průběžné testy + 1 souhrnná práce}\\[0.5mm]
\addlinespace[1.1mm]\multicolumn{2}{@{}l}{\textbf{Září}} & \multicolumn{1}{r@{}}{\textbf{12 hodin}}\\
\midrule
1 & 1.1 & Úvodní hodina: pravidla práce, hodnocení, pomůcky a organizace výuky. \\
2 & 1.1 & Co je výrok? Výroky a nevýroky, pravdivostní hodnota, jednoduchý a složený výrok. \\
3 & 1.1 & Pravdivost složených výroků: negace, konjunkce, disjunkce, implikace, ekvivalence; tabulky. \\
4 & 1.1 & Obecný a existenční kvantifikátor, otevřená formule, negace kvantifikovaných výroků. \\
5 & 1.2 & Číselné obory: přirozená a celá čísla, znaménka, přednost operací. \\
6 & 1.2 & Racionální čísla: zlomky, rozšiřování, krácení a porovnávání. \\
7 & 1.2 & Reálná čísla: iracionální čísla, číselná osa, zaokrouhlování. \\
8 & 1.2 & Intervaly: krajní body, otevřený a uzavřený interval. \\
9 & 1.2 & Absolutní hodnota čísla jako vzdálenost na ose; jednoduché úlohy. \\
10 & T1 & \testr{Průběžný test 1 (15–20 min): Výroková logika a číselné obory. Zbytek hodiny: krátký rozbor a navazující procvičování.} \\
11 & 1.2 & Číselné obory a výroky: pravdivost podle zvolené množiny; oprava chyb z testu. \\
12 & 1.1–1.2 & Množina a její prvky: určení výčtem a vlastností; prázdná množina. \\
\addlinespace[1.1mm]\multicolumn{2}{@{}l}{\textbf{Říjen}} & \multicolumn{1}{r@{}}{\textbf{16 hodin}}\\
\midrule
13 & 1.1–1.2 & Vztahy mezi množinami: podmnožina a rovnost množin. \\
14 & 1.1–1.2 & Vennovy diagramy: znázornění dvou a tří množin. \\
15 & 1.1–1.2 & Sjednocení a průnik množin: význam a zápis. \\
16 & 1.1–1.2 & Rozdíl a doplněk množiny vzhledem k základní množině. \\
17 & 1.1–1.2 & Množinové operace a jejich vlastnosti; kontrola Vennovým diagramem. \\
18 & 1.1–1.2 & Číselné množiny a intervaly: sjednocení a průnik intervalů. \\
19 & 1.1–1.2 & Slovní úlohy na množiny: třídění prvků a počítání počtů. \\
20 & 1.1–1.2 & Opakování množin: zápisy, diagramy, kombinované úlohy. \\
21 & T2 & \testr{Průběžný test 2 (15–20 min): Množiny, intervaly a množinové operace. Zbytek hodiny: krátký rozbor a navazující procvičování.} \\
22 & 1.3 & Pravoúhlý trojúhelník: Pythagorova věta, výpočty délek a úhlů. \\
23 & 1.3 & Sinus, kosinus a tangens ostrého úhlu; poměry stran v pravoúhlém trojúhelníku. \\
24 & 1.3 & Výpočty výšek a vzdáleností pomocí goniometrie; aplikace a kontrola jednotek. \\
25 & 2.1–2.4 & Mocnina s přirozeným exponentem: význam, znaménko, závorky. \\
26 & 2.1–2.4 & Pravidla pro součin a podíl mocnin stejného základu. \\
27 & 2.1–2.4 & Mocnina mocniny, součinu a podílu; odhalování chybných pravidel. \\
28 & 2.1–2.4 & Nulový a záporný celočíselný exponent; převrácená hodnota. \\
\addlinespace[1.1mm]\multicolumn{2}{@{}l}{\textbf{Listopad}} & \multicolumn{1}{r@{}}{\textbf{6 hodin}}\\
\midrule
29 & 2.1–2.4 & Vědecký zápis čísel a řádové odhady. \\
30 & 2.1–2.4 & Druhá a třetí odmocnina; vztah k druhé a třetí mocnině. \\
31 & 2.1–2.4 & Mocniny s racionálním exponentem; převody na odmocniny. \\
32 & P1 & \testr{Čtvrtletní písemná práce: Výroková logika, množiny, goniometrie a mocniny.} \\
33 & 2.1–2.4 & Základní věty pro odmocniny včetně podmínek platnosti. \\
34 & 2.1–2.4 & Vynášení a částečné odmocňování, slučování odmocnin. \\
\addlinespace[1mm]\multicolumn{3}{@{}l}{\footnotesize\textbf{2. čtvrtletí} -- 2 průběžné testy + 1 souhrnná práce}\\[0.5mm]
35 & 2.1–2.4 & Usměrňování jmenovatele s jedním odmocninovým členem. \\
36 & 2.1–2.4 & Usměrňování jmenovatele se sdruženým výrazem. \\
37 & 2.1–2.4 & Složené výrazy s mocninami a odmocninami – postup a podmínky. \\
38 & 2.1–2.4 & Slovní a technické úlohy s mocninami a odmocninami. \\
39 & 2.1–2.4 & Opakování mocnin a odmocnin; cílené procvičení před testem. \\
40 & T3 & \testr{Průběžný test 3 (15–20 min): Mocniny a odmocniny. Zbytek hodiny: krátký rozbor a navazující procvičování.} \\
41 & 3.1–3.5 & Algebraický výraz, proměnná, koeficient, člen a stupeň mnohočlenu. \\
42 & 3.1–3.5 & Sčítání a odčítání mnohočlenů; roznásobování závorek. \\
43 & 3.1–3.5 & Násobení jednočlenu mnohočlenem; znaménka a pořadí úprav. \\
44 & 3.1–3.5 & Násobení dvou mnohočlenů; uspořádání výsledku. \\
\addlinespace[1.1mm]\multicolumn{2}{@{}l}{\textbf{Prosinec}} & \multicolumn{1}{r@{}}{\textbf{12 hodin}}\\
\midrule
45 & 3.1–3.5 & Dělení mnohočlenu jednočlenem; opakování násobení. \\
46 & 3.1–3.5 & Dělení mnohočlenů mnohočlenem: princip algoritmu. \\
47 & 3.1–3.5 & Dělení mnohočlenů se zbytkem; kontrola násobením. \\
48 & 3.1–3.5 & Rozklad mnohočlenu vytýkáním společného činitele. \\
49 & 3.1–3.5 & Vzorce pro druhou mocninu dvojčlenu a rozdíl čtverců. \\
50 & 3.1–3.5 & Kombinace vytýkání a vzorců při rozkladu. \\
51 & 3.1–3.5 & Součet a rozdíl třetích mocnin; vhodná volba vzorce. \\
52 & 3.1–3.5 & Lomené algebraické výrazy: podmínky, krácení. \\
53 & 3.1–3.5 & Sčítání a odčítání lomených algebraických výrazů. \\
54 & 3.1–3.5 & Násobení a dělení lomených výrazů; součinové rozklady. \\
55 & 3.1–3.5 & Složené lomené výrazy; společné podmínky. \\
56 & 3.1–3.5 & Smíšené úpravy polynomů a lomených výrazů. \\
\addlinespace[1.1mm]\multicolumn{2}{@{}l}{\textbf{Leden}} & \multicolumn{1}{r@{}}{\textbf{16 hodin}}\\
\midrule
57 & T4 & \testr{Průběžný test 4 (15–20 min): Algebraické výrazy a rozklady. Zbytek hodiny: krátký rozbor a navazující procvičování.} \\
58 & 3.1–3.5 & Opakování algebraických výrazů; důkaz správnosti dosazením. \\
59 & 4.1–4.4 & Funkce jako přiřazení: argument a funkční hodnota, předpis. \\
60 & 4.1–4.4 & Definiční obor a obor hodnot z předpisu a grafu. \\
61 & 4.1–4.4 & Graf funkce: intervaly růstu a poklesu, průsečíky s osami. \\
62 & 4.1–4.4 & Základní vlastnosti funkcí: monotonie, sudost, lichost. \\
63 & 4.1–4.4 & Lineární funkce a její předpis; graf jako přímka. \\
64 & 4.1–4.4 & Směrnice a úsek na ose: geometrický význam parametrů. \\
65 & 4.1–4.4 & Sestrojení grafu lineární funkce z předpisu a tabulky. \\
66 & 4.1–4.4 & Konstantní funkce; vztah k lineární funkci. \\
67 & 4.1–4.4 & Lineární funkce v technických a ekonomických úlohách. \\
68 & 4.1–4.4 & Grafické porovnávání dvou lineárních funkcí. \\
69 & P2 & \testr{Pololetní písemná práce: Algebraické výrazy a základy funkcí.} \\
70 & 4.1–4.4 & Funkce s absolutní hodnotou: graf y=|x| a základní transformace. \\
71 & 4.1–4.4 & Samostatné procvičování grafů a vlastností lineárních funkcí. \\
72 & 4.1–4.4 & Opakování funkcí; přechod od grafu k rovnici. \\
\addlinespace[1mm]\multicolumn{3}{@{}l}{\footnotesize\textbf{3. čtvrtletí} -- 2 průběžné testy + 1 souhrnná práce}\\[0.5mm]
\addlinespace[1.1mm]\multicolumn{2}{@{}l}{\textbf{Únor}} & \multicolumn{1}{r@{}}{\textbf{12 hodin}}\\
\midrule
73 & 4.5–4.11 & Lineární rovnice: ekvivalentní úpravy a zkouška. \\
74 & 4.5–4.11 & Rovnice se závorkami a zlomky; postupné úpravy. \\
75 & 4.5–4.11 & Lineární nerovnice: změna znaménka při násobení záporným číslem. \\
76 & 4.5–4.11 & Rovnice s parametrem: jeden kořen, žádný kořen, nekonečně mnoho. \\
77 & 4.5–4.11 & Soustava dvou lineárních rovnic: dosazovací metoda. \\
78 & 4.5–4.11 & Soustava dvou lineárních rovnic: sčítací metoda. \\
79 & 4.5–4.11 & Soustavy rovnic: geometrická interpretace a počet řešení. \\
80 & 4.5–4.11 & Rovnice s absolutní hodnotou: vzdálenost na číselné ose. \\
81 & 4.5–4.11 & Rovnice s absolutní hodnotou: rozdělení na případy. \\
82 & 4.5–4.11 & Soustavy lineárních rovnic se třemi neznámými. \\
83 & 4.5–4.11 & Soustavy rovnic: slovní úlohy a kontrola dosazením. \\
84 & T5 & \testr{Průběžný test 5 (15–20 min): Lineární rovnice a nerovnice. Zbytek hodiny: krátký rozbor a navazující procvičování.} \\
\addlinespace[1.1mm]\multicolumn{2}{@{}l}{\textbf{Březen}} & \multicolumn{1}{r@{}}{\textbf{16 hodin}}\\
\midrule
85 & 4.5–4.11 & Smíšené lineární rovnice s parametrem: rozbor případů. \\
86 & 4.5–4.11 & Lineární nerovnice se zlomky a závorkami. \\
87 & 4.5–4.11 & Soustavy lineárních nerovnic: průnik množin řešení. \\
88 & 4.5–4.11 & Intervalový zápis řešení nerovnic a soustav. \\
89 & 4.5–4.11 & Nerovnice s absolutní hodnotou: geometrický význam. \\
90 & 4.5–4.11 & Nerovnice s absolutní hodnotou: rozklad podle znaménka. \\
91 & 4.5–4.11 & Rovnice se dvěma absolutními hodnotami: kritické body. \\
92 & 4.5–4.11 & Nerovnice se dvěma absolutními hodnotami: jednoduché úlohy. \\
93 & 4.5–4.11 & Lineární rovnice a nerovnice: kombinované procvičení. \\
94 & 4.5–4.11 & Slovní úlohy vedoucí na lineární rovnice. \\
95 & 4.5–4.11 & Slovní úlohy vedoucí na soustavy rovnic. \\
96 & 4.5–4.11 & Vyjádření neznámé z technického vzorce a návaznost na rovnice. \\
97 & 4.5–4.11 & Vyjádření neznámé z lomeného technického vzorce; podmínky. \\
98 & T6 & \testr{Průběžný test 6 (15–20 min): Soustavy a rovnice s absolutní hodnotou. Zbytek hodiny: krátký rozbor a navazující procvičování.} \\
99 & 4.5–4.11 & Procvičení soustav rovnic a nerovnic – volba metody. \\
100 & 4.5–4.11 & Opakování lineárních rovnic s absolutní hodnotou. \\
\addlinespace[1.1mm]\multicolumn{2}{@{}l}{\textbf{Duben}} & \multicolumn{1}{r@{}}{\textbf{12 hodin}}\\
\midrule
101 & 4.5–4.11 & Diagnostika typických chyb a oprava předchozích úloh. \\
102 & 4.5–4.11 & Souhrnné opakování lineárních rovnic, nerovnic a soustav. \\
103 & 5.1–5.9 & Kvadratická funkce: parabola, definiční obor a obor hodnot. \\
104 & 5.1–5.9 & Graf kvadratické funkce y=ax²; vliv parametru a. \\
105 & 5.1–5.9 & Posuny paraboly a vrcholový tvar; souvislost s osou souměrnosti. \\
106 & 5.1–5.9 & Doplnění na čtverec: získání vrcholového tvaru. \\
107 & 5.1–5.9 & Kvadratický trojčlen a průsečíky paraboly s osami. \\
108 & 5.1–5.9 & Kvadratická funkce v obecném tvaru: konstrukce grafu. \\
109 & 5.1–5.9 & Grafy kvadratických funkcí – samostatná práce. \\
110 & P3 & \testr{Čtvrtletní písemná práce: Lineární funkce a rovnice, úvod kvadratických funkcí.} \\
111 & 5.1–5.9 & Neúplné kvadratické rovnice: vytýkání a odmocnění. \\
112 & 5.1–5.9 & Úplná kvadratická rovnice: diskriminant a vzorec. \\
\addlinespace[1mm]\multicolumn{3}{@{}l}{\footnotesize\textbf{4. čtvrtletí} -- 2 průběžné testy + 1 souhrnná práce}\\[0.5mm]
\addlinespace[1.1mm]\multicolumn{2}{@{}l}{\textbf{Květen}} & \multicolumn{1}{r@{}}{\textbf{12 hodin}}\\
\midrule
113 & 5.1–5.9 & Kvadratická rovnice: rozklad na součin a nulový součin. \\
114 & 5.1–5.9 & Viètovy vztahy: součet a součin kořenů. \\
115 & 5.1–5.9 & Sestavení kvadratické rovnice z daných kořenů. \\
116 & 5.1–5.9 & Grafické a algebraické řešení kvadratických rovnic – srovnání. \\
117 & 5.1–5.9 & Rovnice vedoucí na kvadratickou rovnici: jednoduchá substituce. \\
118 & 5.1–5.9 & Soustava lineární a kvadratické rovnice o dvou neznámých. \\
119 & 5.1–5.9 & Soustavy rovnic s parabolou: geometrický význam průsečíků. \\
120 & 5.1–5.9 & Rovnice s odmocninou: podmínky a ověřování kořenů. \\
121 & T7 & \testr{Průběžný test 7 (15–20 min): Kvadratické rovnice a grafy. Zbytek hodiny: krátký rozbor a navazující procvičování.} \\
122 & 5.1–5.9 & Rovnice se dvěma odmocninami: izolace a kontrola. \\
123 & 5.1–5.9 & Kvadratická rovnice s parametrem: zvláštní případy. \\
124 & 5.1–5.9 & Kvadratická nerovnice: odečítání řešení z grafu paraboly. \\
\addlinespace[1.1mm]\multicolumn{2}{@{}l}{\textbf{Červen}} & \multicolumn{1}{r@{}}{\textbf{12 hodin}}\\
\midrule
125 & 5.1–5.9 & Kvadratická nerovnice: metoda nulových bodů a znamének. \\
126 & 5.1–5.9 & Nerovnice s kvadratickým trojčlenem a intervaly. \\
127 & 5.1–5.9 & Souhrnné úlohy na kvadratické rovnice a nerovnice. \\
128 & 5.1–5.9 & Slovní úlohy vedoucí na kvadratickou rovnici. \\
129 & 5.1–5.9 & Grafická optimalizace jednoduché kvadratické funkce. \\
130 & T8 & \testr{Průběžný test 8 (15–20 min): Kvadratické nerovnice a odmocninové rovnice. Zbytek hodiny: krátký rozbor a navazující procvičování.} \\
131 & 5.1–5.9 & Technické aplikace: obsah a rozměry, kontrola fyzikální smysluplnosti. \\
132 & 5.1–5.9 & Opakování kvadratických funkcí: graf, vrchol, nulové body. \\
133 & 5.1–5.9 & Opakování kvadratických rovnic: metody a volba postupu. \\
134 & 5.1–5.9 & Oprava typických chyb z kvadratických nerovnic. \\
135 & P4 & \testr{Pololetní písemná práce: Souhrnné učivo 1. ročníku.} \\
136 & 5.1–5.9 & Závěrečné procvičení a shrnutí prvního ročníku. \\
\end{longtable}
\normalsize
\end{document}
